\appendixitem{REPRESENTATIONS OF REAL ALGEBRAS}

Denote by \(\sigma\) the automorphism \(\alpha\mapsto\bar{\alpha}\) of C; if W is a complex vector space, denote by \(\overline{W}\) the C-vector space \(\sigma_{*}(W)\) (that is, the group W with the law of operation \((\alpha,w)\mapsto\bar{\alpha}w\) for \(\alpha\in C, w\in W\) ).

PROPOSITION 1. Let A be an R-algebra (associative and unital) and V a finite dimensional simple A-module over R. Then, we must be in one of the following three situations:

\(\alpha)\) The commutant of V (Algebra, Chap. VIII, § 5, no. 1) is isomorphic to \(\mathbf{R}\) , and the \(\mathrm{A}_{(\mathbf{C})}\) -module \(\mathrm{V}_{(\mathbf{C})}\) is simple;

\(\beta)\) the commutant of V is isomorphic to \(\mathbf{C}\) ; the \(\mathrm{A}_{(\mathbf{C})}\) -module \(\mathrm{V}_{(\mathbf{C})}\) is the direct sum of two non-isomorphic simple \(\mathrm{A}_{(\mathbf{C})}\) -submodules which are interchanged by \(\sigma \otimes 1_{\mathrm{V}}\) ;

\(\gamma)\) the commutant of V is isomorphic to H; the \(A_{(C)}\) -module \(V_{(C)}\) is the direct sum of two isomorphic simple \(A_{(C)}\) -submodules that are interchanged by \(\sigma \otimes 1_V\) .

The commutant E of V is a field, a finite extension of R (Algebra, Chap. VIII, §3, no. 2, Prop. 2), hence isomorphic to R, C or H (Algebra, Chap. VIII, §15). The \(A_{(C)}\) -module \(V_{(C)}\) is semi-simple (Algebra, Chap. VIII, §11, no. 4), and its commutant can be identified with \(C \otimes_{R} E\) (Algebra, Chap. VIII, §11, no. 2, Lemma 1).

If \(\mathbf{E}\) is isomorphic to \(\mathbf{R}\) , the commutant of \(V_{(\mathbf{C})}\) is isomorphic to \(\mathbf{C}\) , and \(V_{(\mathbf{C})}\) is a simple \(A_{(\mathbf{C})}\) -module (Algebra, Chap. VIII, §11, no. 4).

If \(\mathbf{E}\) is not isomorphic to \(\mathbf{R}\) , it contains a field isomorphic to \(\mathbf{C}\) ; it follows that \(V\) has an \(A_{(\mathbf{C})}\) -module structure, denoted by \(V^c\) . Then \(V^c\) is a simple

\(A_{(C)}\) -module, and the C-linear map \(\psi: V_{(C)} \to V^{c} \oplus \overline{V}^{c}\) such that \(\psi(\alpha \otimes v) = (\alpha v, \bar{\alpha} v)\) for \(\alpha \in C\) , \(v \in V\) , is an isomorphism (Algebra, Chap. V, §10, no. 4, Prop. 8). Moreover, \(\sigma \otimes 1_{V}\) corresponds under this isomorphism to the R-automorphism \((v, v') \mapsto (v', v)\) of \(V^{c} \oplus \overline{V}^{c}\) , and hence interchanges the two \(A_{(C)}\) -submodules \(\psi^{-1}(V^{c})\) and \(\psi^{-1}(\overline{V}^{c})\) .

The commutant \(E_{(\mathbf{C})}\) of \(V_{(\mathbf{C})}\) thus contains \(C \times C\) , operating by homotheties on \(V^{c} \oplus \overline{V}^{c}\) . There is no isomorphism of \(A_{(\mathbf{C})}\) -modules from \(V^{c}\) to \(\overline{V}^{c}\) if and only if \(E_{(\mathbf{C})}\) reduces to \(C \times C\) , that is, if E is isomorphic to C. This completes the proof.

PROPOSITION 2. Let A be an R-algebra (associative and unital), and W a finite dimensional simple \(\mathrm{A}_{(\mathbf{C})}\) -module over C. Then, we must be in one of the following three situations:

\begin{enumerate}
\def\labelenumi{\alph{enumi})}
\item
  There exists an \(A_{(\mathbf{C})}\) -isomorphism \(\theta\) from W to \(\overline{W}\) with \(\theta \circ \theta = 1_{W}\) . Then the set V of fixed points of \(\theta\) is an R-structure on W, and a simple A-module with commutant R.1 \(_{V}\) . Moreover, \(W_{[R]}\) is the direct sum of two isomorphic simple A-modules.
\item
  The \(\mathrm{A}_{(\mathbf{C})}\) -modules W and \(\overline{\mathrm{W}}\) are not isomorphic; then \(\mathrm{W}_{[\mathbf{R}]}\) is a simple A-module with commutant \(\mathbf{C}.1_{\mathrm{W}}\) .
\item
  There exists an \(\mathrm{A}_{(\mathbf{C})}\) -isomorphism \(\theta\) from W to \(\overline{W}\) with \(\theta \circ \theta = -1_{W}\) . Then the A-module \(W_{[R]}\) is simple, and its commutant is the field \(C.1_{W} \oplus C.\theta\) , which is isomorphic to H.
\end{enumerate}

The complex vector space \(\mathrm{Hom}_{\mathrm{A}_{(\mathbf{C})}}(\mathrm{W},\overline{\mathrm{W}})\) is of dimension \(\leq 1\) (Algebra, Chap. VIII, §3, no. 2); if \(\theta \in \mathrm{Hom}_{\mathrm{A}_{(\mathbf{C})}}(\mathrm{W},\overline{\mathrm{W}})\) , the endomorphism \(\theta \circ \theta\) of W is a homothety, with ratio \(\alpha \in \mathbf{C}\) . For all \(w \in \mathrm{W}\) , we have \(\alpha \theta(w) = \theta \circ \theta \circ \theta(w) = \theta(\alpha w) = \bar{\alpha} \theta(w)\) , so \(\alpha\) is real. If \(\theta' = \lambda \theta\) , with \(\lambda \in \mathbf{C}\) , then \(\theta' \circ \theta' = |\lambda|^2 \theta \circ \theta\) ; thus, exactly one of the following three possibilities is realised:

\begin{enumerate}
\def\labelenumi{\alph{enumi})}
\item
  There exists \(\theta \in \mathrm{Hom}_{\mathrm{A}_{(\mathbf{C})}}(\mathrm{W},\overline{\overline{\mathrm{W}}})\) with \(\theta \circ \theta = 1_{\mathrm{W}}\) ;
\item
  \(\mathrm{Hom}_{\mathrm{A}_{(\mathbf{C})}}(\mathrm{W},\overline{\mathrm{W}}) = \{0\}\) ;
\item
  There exists \(\theta \in \mathrm{Hom}_{\mathrm{A}_{(\mathbf{C})}}(\mathrm{W},\overline{\mathrm{W}})\) with \(\theta \circ \theta = -1_{\mathrm{W}}\) .
\end{enumerate}

In case a), the set V of fixed points of \(\theta\) is an R-structure on W (Algebra, Chap. V, p.~61, Prop. 7); since \(V_{(\mathbf{C})}\) is isomorphic to W, the A-module V is simple with commutant \(R.1_{V}\) (Prop. 1), and \(W_{[R]}\) is not simple.

Conversely, if \(W_{[R]}\) is not simple, let V be a simple A-submodule of \(W_{[R]}\) ; since the \(A_{(C)}\) -module W is simple, \(V + iV = W\) and \(V \cap iV = \{0\}\) , that is, \(W = V \oplus iV\) . Thus, V is an R-structure on W, and the isomorphism \(\theta\) from W to \(\overline{W}\) such that \(\theta(v + iv') = v - iv'\) for v and \(v'\) in V satisfies \(\theta \circ \theta = 1_{W}\) .

Consequently, in cases \(b)\) and \(c)\) , the A-module \(\mathrm{W}_{[\mathbf{R}]}\) is simple; by Prop. 1, its commutant E is isomorphic to C in case \(b)\) , and to H in case \(c)\) . Moreover, it is clear that E contains \(\mathbf{C}.1_{\mathrm{W}}\) , and \(\mathbf{C}.\theta\) in case \(c)\) , hence the proposition.

With the assumptions in the proposition, the \(A_{(C)}\) -module W is said to be of real, complex or quaternionic type (relative to A) in case a), b) or c), respectively.

For K = R or C, denote by \(\mathfrak{S}_{\mathrm{K}}(\mathrm{A})\) the set of classes of finite dimensional simple \(A_{(K)}\) -modules over K. The group \(\Gamma = \operatorname{Gal}(\mathbf{C}/\mathbf{R})\) operates on \(\mathfrak{S}_{\mathbf{C}}(\mathrm{A})\) ; the two preceding propositions establish a bijective correspondence between \(\mathfrak{S}_{\mathbf{R}}(\mathrm{A})\) and the quotient set \(\mathfrak{S}_{\mathbf{C}}(\mathrm{A})/\Gamma\) .
% semantic-fragments: legacy-input-seam
\relax{}
